\chapterrule
\section{从可调函数到损失：怎样比较预测？}
\label{intro:00-functions}
\label{intro:16-loss-gradient}

序章里，“下雨”“出门”“一把”让“伞”成为一个合理的后续候选。
换一种表达或上下文，合适的后续内容也可能变化。
想靠人逐条写出各种表达与语境下的预测规则，很困难。
另一个办法是选\textbf{带参数的函数}，让训练样本参与决定参数。

为了先看清参数怎样由样本确定，我们先用一个独立的数值例子练习，
后面再把这套学习思路用于文字预测。
本章与下一章始终使用三个样本 $(-1,-1),(0,2),(1,2)$：
先认识参数，再用损失比较预测；下一章学习怎样逐步改善参数。


\subsection{函数是一台按规则输出结果的机器}
\phantomsection\label{concept:function}

把一个数 $x$ 送入函数 $f$，得到输出 $y=f(x)$。例如
\[
f(x)=2x+1,\qquad f(0)=1,\qquad f(1)=3.
\]
输入也可以是一组数字。预测后续文字时，可以先把前文表示成一组数字，
再用函数为后续候选计算分数；这些步骤将在后面的章节逐步展开。
函数给我们提供了一个可计算、可检验的预测模型。

以前我们会先知道公式，再计算结果；机器学习常常反过来：先看到一些输入与结果，再寻找一个能够解释这些例子的函数。

\subsection{三个样本，一条可调的直线}
\phantomsection\label{concept:sample}

设训练数据是
\[
(-1,-1),\quad(0,2),\quad(1,2).
\]
我们先选一个最简单的函数族：
\begin{equation}
\widehat{y}=wx+b.
\label{eq:line-first}
\end{equation}
$\widehat{y}$ 表示模型的\textbf{预测值}，而 $y$ 是实际观察到的值。$w$ 决定斜率，$b$ 决定截距。

先令 $w=1,b=0$，模型成为 $\widehat{y}=x$。
三个输入依次得到 $(-1,0,1)$，与真实值 $(-1,2,2)$ 比较，残差为 $(0,-2,-1)$。
只有第一点预测正确。保持 $w=1$，把 $b$ 从 $0$ 调到 $1$，预测变成
\[
\widehat{y}(-1)=0,\qquad \widehat{y}(0)=1,\qquad \widehat{y}(1)=2.
\]
第二、三个样本的预测更接近真实值了，第一个样本却从猜对变成了猜错。
若再把 $b$ 加到 $2$，中间的样本猜对了，两端的预测又都偏高。
参数调整会同时影响多个样本，需要比较整组预测的表现。

为什么两个参数各有用处？增大 $b$ 会把三个预测一起向上移；
增大 $w$ 则让 $x=-1$ 的预测下降、$x=1$ 的预测上升，$x=0$ 的预测保持不变。
这三个点也无法被同一条直线全部穿过：输入每次增加 $1$，
输出却先增加 $3$、再增加 $0$，而直线每次的输出增量相同。
我们要寻找一组能让整体预测更合适的参数。

\par\noindent\begin{minipage}{\linewidth}
\tagstructbegin{tag=Figure,alttext={三个不共线的观察点负一负一、零二、一二。三条候选直线分别为 y 等于 x、x 加一、x 加二；增大偏置使直线整体上移，某个点更准时另一个点可能变差，没有一条直线穿过全部三点。}}
\begin{center}
\begin{tikzpicture}
\begin{axis}[
width=.77\linewidth,height=57mm,
axis lines=middle,xmin=-1.7,xmax=1.7,ymin=-2.2,ymax=4.2,
xlabel={$x$},ylabel={$y$},grid=both,
grid style={linegray,thin},legend pos=north west,
legend style={font=\footnotesize},
xtick={-1,0,1},ytick={-1,0,1,2,3},
]
\addplot[datacolor,thick,dashed,domain=-1.4:1.35,samples=2]{x};
\addlegendentry{候选 A：$\hat{y}=x$}
\addplot[datacolor,very thick,domain=-1.4:1.35,samples=2]{x+1};
\addlegendentry{候选 B：$\hat{y}=x+1$}
\addplot[datacolor,thick,dashdotted,domain=-1.4:1.35,samples=2]{x+2};
\addlegendentry{候选 C：$\hat{y}=x+2$}
\addplot[only marks,mark=*,noisecolor,mark size=2.5pt] coordinates {(-1,-1) (0,2) (1,2)};
\addlegendentry{观察数据}
\end{axis}
\end{tikzpicture}
\captionof{figure}{同一组三点与三条候选直线：调整偏置会同时改变各个样本的预测}
\label{fig:first-line-fit}
\end{center}
\tagstructend
\end{minipage}\par
\medskip

\subsection{从可调函数到模型、训练与推理}
\phantomsection\label{concept:model}
\phantomsection\label{concept:parameter}
\phantomsection\label{concept:weight}
\phantomsection\label{concept:training}
\phantomsection\label{concept:inference}

请先分清四个词：
\begin{enumerate}
\item \textbf{模型（model）}：我们选择的函数形式，例如 $\widehat{y}=wx+b$。
\item \textbf{参数（parameters）}：模型内部可调的数，例如 $w,b$。
\item \textbf{训练（training）}：根据样本调整参数，让预测更符合任务目标。
\item \textbf{推理（inference）}：参数确定后，输入新数据，计算预测结果。
\end{enumerate}
在深度学习里，函数可以有数百万乃至更多的可训练参数；但计算过程仍是\textbf{输入、运算、输出}。

拟合需要用新数据检验：模型在已有样本上表现好，对新数据仍可能预测不准。后续的\textbf{泛化}章节专门讨论这个问题。
